All methods share one nnU-Net~\cite{isensee2021nnunet} configuration; only
the input augmentation differs. Nothing below is tuned per dataset.

\paragraph{Training schedule.} 3D full-resolution nnU-Net, SGD with Nesterov
momentum $0.99$, initial learning rate $10^{-2}$ with the nnU-Net
\texttt{polyLR} decay, $250$ mini-batches per epoch, combined Dice and
cross-entropy loss, deep supervision on. Patch size, spacing and batch size
are whatever the nnU-Net planner selects for that dataset, identically for
every method. The epoch budget is fixed per dataset and shared by all methods
on it: $2500$ (Brats-GLI), $2000$ (ON-Harmony), $2000$ (Open-MS), $200$
(CHAOS). Every number reported in this paper is a 3-fold cross-validated
average; folds are the same across methods.

\paragraph{PALETTE hyper-parameters.} Set once and never tuned per dataset:
$k$-means class count drawn per sample from $\{2,\dots,6\}$, skipped in
favour of a single global remap with probability $0.10$; Voronoi
sub-parcellation of each intensity class with probability $0.60$ (else left
as one region); signed affine coefficient $\alpha \in \pm U(0.5, 2.0)$
applied per sub-region as in \cref{eq:affine} (always applied, to whatever
partition results from the $k$-means/Voronoi steps); per-label affine
decoupling applied independently per label with probability $0.5$ (skipped
for any label with under $4$ foreground voxels in that sample); Gaussian
blur $\sigma$ drawn from $\{0,0,0,0.3,0.5,0.8\}$, applied twice (after the
region-level and after the label-level remap), so mostly not applied at
either point. The transform is applied on the GPU, on-the-fly, to a fixed
$50\%$ of training
iterations (\emph{train050}); the remaining half of iterations use the
un-transformed real image.

\paragraph{Validation-time configuration.} \emph{val000} (the configuration
reported in the main paper) applies no synthetic contrast to the validation
split, matching every competing method and keeping model selection
like-for-like. \emph{val100} applies PALETTE to the entire validation split;
it is reported in \cref{sec:suppl-val100}.

\paragraph{Baseline configurations.} SynthSeg-noEM uses the authors' released
generator with its default priors. SynthSeg-EM is our re-implementation of
the variant described but not released in the original work; it estimates a
per-class intensity model from the sparse task labels available in each
dataset via expectation-maximization, and is otherwise identical in its
sampling procedure. SRCSM and Auglab use their published default
configurations. No competing method was re-tuned in our favour, and all use
the same folds, schedule and epoch budget as above.

\paragraph{Compute.} Training was performed on NVIDIA L40S and H100 GPUs. A
single fold costs from under an hour (CHAOS, $200$ epochs) to roughly three
days (Brats-GLI, $2500$ epochs); the full study is dominated by the $7$
methods $\times$ $8$ settings $\times$ $3$ folds main comparison plus the
ablation ladders of \cref{subsec:e-ablation}.
